\section{A Self-Caught Audit Error}
\label{sec:retraction}

Between the characterisation and the intervention we planned a fourth
module, A, to fix the carry-aware overlap check. Our own Phase 3 audit had
reported that the check false-accepts $55.4\%$ of all possible overlap bit
patterns --- serious over-permissiveness in a mechanism central to the
published method, and consistent with the $43.8\%$ success rate under
overlap corruption in Sec.~\ref{sec:adversarial}.

Re-validating that evidence before writing any code showed the finding was
wrong, and wrong in the audit rather than in the mechanism. \textbf{The audit
script pinned the actual computed overlap to one bit regardless of the
nominal overlap width it was sweeping}, so the overlapping bit strings were
never compared at their declared width. The swept parameter did not vary the
quantity it named.

Three independent checks establish the correction.

\begin{enumerate}\itemsep2pt
\item \emph{Algebraically}, the implementation's acceptance condition and the
audit's ground-truth oracle reduce to the same Boolean function, so the
original test could only ever have detected its own bug. The companion
article~\cite{methodpaper} states and proves this.
\item \emph{Empirically}, the same exhaustive 168-combination audit re-run on
corrected geometry gives \textbf{0\% false-accept and 0\% false-reject}
across all three overlap widths --- 42 true accepts and 126 true rejects,
with no error of either kind (\SIaudit{}).
\item \emph{By instrumented replay}, a deliberately corrupted candidate was
rejected in all 80 trials of the overlap-corruption condition.
\end{enumerate}

Two consequences follow. Module A was deleted from the design before any
code was written, and the carry-check implementation required no change. The
residual failures under overlap corruption are explained by the same
budget-starvation mechanism as everything else in Sec.~\ref{sec:attribution}
--- the corrupted candidate is correctly rejected, and a narrow beam then has
no alternative path left --- rather than by the check being fooled.

We report this at length for two reasons. The corrected result is what
justifies leaving the carry check untouched, so the reader needs it to
evaluate the intervention in Part II. And the failure pattern it exposes ---
a swept parameter that does not in fact vary the quantity it names --- is a
reusable warning for anyone building comparable audits. The prior reports
were corrected in place with visible strikethrough rather than silently
edited.

% =====================================================================
